\clearpage

\subsection{Диапазон наблюдений, активация и восстановление}

Таблица при нагрузочных возвратах показывает выбор retained-load. Числа порогов обозначают фактические значения WIP и число обучений. Диапазон обучения — разброс максимального наблюдавшегося WIP между наборами; $q/\max w$ считается внутри каждого набора. K/D обозначают известную/обнаруженную структуру. Диагностика не меняет правила активации.

\begin{center}\small\setlength{\tabcolsep}{3pt}\begin{tabular}{llp{.24\textwidth}rrr}\toprule

Process & N & Threshold:count & Train max & $q/\max w$ & Above\\\midrule

P1 & K & 6:9; 12:3 & 5--9 & 0.67--2.00 & 6/12\\

P1 & D & 6:9; 12:3 & 5--9 & 0.67--2.00 & 6/12\\

P2 & K & 8:9; 16:3 & 6--13 & 0.62--2.67 & 6/12\\

P2 & D & 8:9; 16:3 & 6--13 & 0.62--2.67 & 6/12\\

P3 & K & 10:11; 40:1 & 7--11 & 0.91--5.00 & 7/12\\

P3 & D & 10:11; 40:1 & 7--11 & 0.91--5.00 & 7/12\\

P4 & K & 10:11; 20:1 & 10--21 & 0.48--1.43 & 1/12\\

P4 & D & 10:11; 20:1 & 10--21 & 0.48--1.43 & 1/12\\

P5 & K & 10:12 & 11--18 & 0.56--0.91 & 0/12\\

P5 & D & 10:12 & 11--18 & 0.56--0.91 & 0/12\\

P6 & K & 10:11; 20:1 & 10--16 & 0.62--2.00 & 1/12\\

P6 & D & 10:11; 20:1 & 10--16 & 0.62--2.00 & 1/12\\

\bottomrule\end{tabular}\end{center}

\begin{center}\small\setlength{\tabcolsep}{4pt}\begin{tabular}{llrrrrr}\toprule

Process & N & Activate & Nonsevere & Lead/$T$ & Event & Prevented\\\midrule

P1 & K & 0.705 & 0.342 & 5.901 & 0.365 & 0.329\\

P1 & D & 0.705 & 0.342 & 5.901 & 0.365 & 0.329\\

P2 & K & 0.705 & 0.490 & 5.042 & 0.345 & 0.179\\

P2 & D & 0.705 & 0.490 & 5.042 & 0.345 & 0.179\\

P3 & K & 0.677 & 0.031 & 4.247 & 0.642 & 0.036\\

P3 & D & 0.677 & 0.031 & 4.247 & 0.642 & 0.036\\

P4 & K & 0.812 & 0.506 & 7.837 & 0.495 & 0.200\\

P4 & D & 0.812 & 0.506 & 7.837 & 0.495 & 0.200\\

P5 & K & 0.868 & 0.720 & 9.730 & 0.382 & 0.274\\

P5 & D & 0.868 & 0.720 & 9.730 & 0.382 & 0.274\\

P6 & K & 0.762 & 0.320 & 6.301 & 0.382 & 0.412\\

P6 & D & 0.762 & 0.320 & 6.301 & 0.382 & 0.412\\

\bottomrule\end{tabular}\end{center}

Nonsevere — активация при отсутствии тяжёлого события в парной траектории без действия; это не обязательно вредное вмешательство. Время предупреждения измеряется относительно начала такого контрфактического события. Предотвращение условно по событию без действия. Восстановление ниже означает WIP не выше $2S$ в течение $10T$ после импульса с цензурированием на исходном горизонте $180T$. Сначала усредняются траектории каждого независимого обучения, затем сами обучения; учитываются все импульсные траектории, а не только кризисные.

\begin{center}\small\begin{tabular}{lrrrr}\toprule

Process & No action & Predictive K & Retained K & Generator\\\midrule

P1 & 0.771 & 1.000 & 1.000 & 1.000\\

P2 & 0.745 & 0.979 & 0.979 & 1.000\\

P3 & 0.195 & 0.659 & 0.599 & 0.732\\

P4 & 0.536 & 0.982 & 0.867 & 0.987\\

P5 & 0.802 & 0.865 & 0.943 & 1.000\\

P6 & 0.482 & 0.948 & 0.987 & 1.000\\

\bottomrule\end{tabular}\end{center}

\subsection{Последствия общих ресурсов и положение узкого места}

Ресурсная диагностика повторяет P4/P5 при всех четырёх механизмах и всех 12 обучениях для двух импульсных условий с высокой базовой интенсивностью и повтора 0. Включены все типы действий на минимальном фиксированном пороге и сохранённые выбранные политики. Все 884 префиксов до $180T$ точно совпадают с исходными тестовыми отсчётами. Затем поступления прекращаются, а действия сохраняются до $600T$ для завершения когорты. Интегралы очередей относятся только к исходному горизонту. Поэтому длительность полной когорты имеет другой горизонт, чем основной показатель.

\begin{longtable}{llrrrr}\caption{P4: интегралы WIP каждого процесса (тысячи) по полному тесту и времена завершения полной когорты по продлённой диагностике.}\\\toprule

Mechanism & Action at 10 & WIP X & WIP Y & Cycle X & Cycle Y\\\midrule\endfirsthead

\toprule Mechanism & Action at 10 & WIP X & WIP Y & Cycle X & Cycle Y\\\midrule\endhead

fixed & None & 7.64 & 8.41 & 23.27 & 25.61\\

fixed & Capacity & 4.62 & 5.31 & 13.21 & 15.03\\

fixed & Rework X & 7.33 & 8.26 & 22.19 & 25.05\\

fixed & Rework Y & 7.34 & 7.82 & 22.09 & 23.40\\

service & None & 43.69 & 45.71 & 176.41 & 181.78\\

service & Capacity & 22.70 & 24.06 & 75.40 & 78.33\\

service & Rework X & 41.75 & 44.77 & 165.43 & 175.66\\

service & Rework Y & 42.15 & 42.57 & 166.09 & 164.01\\

routing & None & 31.65 & 30.88 & 119.63 & 116.00\\

routing & Capacity & 15.28 & 15.16 & 51.39 & 50.36\\

routing & Rework X & 16.83 & 22.30 & 58.65 & 77.74\\

routing & Rework Y & 21.92 & 16.61 & 76.21 & 56.21\\

combined & None & 90.99 & 89.23 & 538.82 & 525.69\\

combined & Capacity & 71.85 & 70.13 & 347.75 & 338.59\\

combined & Rework X & 69.00 & 83.06 & 333.96 & 441.36\\

combined & Rework Y & 85.34 & 69.00 & 456.45 & 335.60\\

\bottomrule\end{longtable}

\begin{longtable}{llrrr}\caption{P5: средние очереди B и D на исходном горизонте; число парных траекторий с перемещением наибольшей очереди от B к D.}\\\toprule

Mechanism & Action at 10 & Queue B & Queue D & B to D\\\midrule\endfirsthead

\toprule Mechanism & Action at 10 & Queue B & Queue D & B to D\\\midrule\endhead

fixed & None & 12.13 & 0.11 & 0/24\\

fixed & Capacity B & 1.76 & 4.55 & 24/24\\

fixed & Capacity D & 12.13 & 0.01 & 0/24\\

fixed & Rework & 12.13 & 0.08 & 0/24\\

service & None & 140.59 & 0.01 & 0/24\\

service & Capacity B & 38.34 & 0.11 & 0/24\\

service & Capacity D & 140.59 & 0.00 & 0/24\\

service & Rework & 140.59 & 0.00 & 0/24\\

routing & None & 12.17 & 36.45 & 0/24\\

routing & Capacity B & 1.67 & 46.10 & 0/24\\

routing & Capacity D & 12.17 & 4.38 & 0/24\\

routing & Rework & 12.17 & 1.73 & 0/24\\

combined & None & 137.95 & 1.28 & 0/24\\

combined & Capacity B & 37.29 & 116.13 & 24/24\\

combined & Capacity D & 137.95 & 0.05 & 0/24\\

combined & Rework & 137.95 & 0.04 & 0/24\\

\bottomrule\end{longtable}
